% ERDŐS #404 — THE p-ADIC FACTORIAL-SUM REDUCTION: COMPLETE DERIVATION
% Round 3 of the bulldoze campaign (2026-11-02).  Companion kernel:
% problem404-full/PadicReduction.lean (core Lean 4, CI-verified).
%
% THE PROBLEM (verbatim, erdosproblems.com/404):
%   For which integers a >= 1 and primes p is there a finite upper bound on
%   those k such that there are a = a_1 < ... < a_n with p^k | (a_1!+...+a_n!)?
%   If f(a,p) is the greatest such k, how does this function behave?
%   Is there a prime p and an infinite sequence a_1 < a_2 < ... such that if
%   p^{m_k} is the highest power of p dividing sum_{i<=k} a_i! then m_k -> inf?
\documentclass{article}
\usepackage{amsmath,amssymb,amsthm}
\newtheorem{theorem}{Theorem}
\newtheorem{lemma}[theorem]{Lemma}
\newtheorem{corollary}[theorem]{Corollary}
\newtheorem{definition}[theorem]{Definition}
\newtheorem{conjecture}[theorem]{Conjecture}
\newcommand{\vpp}{v_p}
\newcommand{\Zp}{\mathbb{Z}_p}

\begin{document}
\title{Erd\H{o}s Problem \#404: the $p$-adic factorial-sum reduction,\\
the $p \mid a{+}1$ class, $f(4,2)=6$, and the unbounded candidates}
\date{2026-11-02}
\maketitle

\section{The reduction (Theorem A)}

Fix $a \ge 1$ and let all index sets $B$ range over finite subsets of
$\{a+1, a+2, \ldots\}$.  Define the \textbf{quotient staircase}
\[
  P_i := \frac{i!}{a!} = (a+1)(a+2)\cdots i, \qquad i \ge a+1 .
\]

\begin{theorem}[Factorization; kernel-green]\label{thm:A}
For every finite $B \subseteq \{a{+}1, a{+}2,\ldots\}$,
\[ a! + \sum_{i\in B} i! \;=\; a! \cdot \Bigl(1 + \sum_{i \in B} P_i\Bigr). \]
Consequently, for every prime $p$,
\[ \vpp\Bigl(a! + \sum_{i\in B} i!\Bigr) \;=\; \vpp(a!) + \vpp\Bigl(1+\sum_{i\in B}P_i\Bigr). \]
\end{theorem}
\begin{proof}
Each $i! = a!\,P_i$; factor $a!$ out of every term.  The valuation identity
is the multiplicativity $\vpp(xy) = \vpp(x)+\vpp(y)$ together with $a! \ge 1$
and $1 + \sum P_i \ge 1$.  (Formalized: \texttt{erdos404\_split}; the $p=2$
valuation machinery including multiplicativity is in the kernel, mirroring
the Kill-\#4 \texttt{KernelLaw.lean}.)
\end{proof}

\begin{definition}
$f(a,p) := \sup \bigl\{ \vpp(a! + \sum_{i\in B} i!) : B \text{ finite} \bigr\}$
and $g(a,p) := \sup \bigl\{ \vpp(1 + \sum_{i\in B}P_i) \bigr\} \ge 0$,
so Theorem~\ref{thm:A} gives
\[ f(a,p) = \vpp(a!) + g(a,p). \]
For $k \ge 1$ let $R_k \subseteq \mathbb{Z}/p^k$ be the set of subset sums of
$\{P_i \bmod p^k : \vpp(P_i) < k\}$ (a \emph{finite} multiset: $\vpp(P_i) \to \infty$).
\end{definition}

\begin{lemma}[Prefix property]\label{lem:prefix}
$-1 \in R_{k+1} \Rightarrow -1 \in R_k$.  Hence $\{k : -1 \in R_k\}$ is a
prefix $\{1,\ldots,g\}$, and $g(a,p) = \max\{k : -1 \in R_k\}$ (or $\infty$).
\end{lemma}
\begin{proof}
Reduce a subset sum that is $\equiv -1 \pmod{p^{k+1}}$ modulo $p^k$.
\end{proof}

\begin{lemma}[Lift lemma]\label{lem:lift}
Let $-1 \in R_k$ with witness $B$, and write $\sum_{j\in B}P_j = -1 + p^k c$.
Suppose there is a \textbf{fresh gadget of exact valuation $k$ with unit part
$-c$}: a finite set $G$ of indices beyond $a$, disjoint from $B$, with
$\sum_{i\in G}P_i = p^k u$ where $u \equiv -c \pmod p$ (automatically
$p \nmid u$ when $p \nmid c$; $u = 0$, i.e.\ $G = \emptyset$, when $p \mid c$).
Then $B' = B \cup G$ witnesses $-1 \in R_{k+1}$.
In particular, if the level-$k$ gadget supply contains unit parts covering
all of $\mathbb{Z}/p$, then $-1 \in R_{k+1}$.
\end{lemma}
\begin{proof}
$\sum_{j \in B'}P_j = -1 + p^k(c + u) \equiv -1 \pmod{p^{k+1}}$ since
$p \mid c + u$.
\end{proof}

\begin{theorem}[Compactness / the chain]\label{thm:chain}
The following are equivalent:
\begin{enumerate}
\item $g(a,p) = \infty$ (no level obstructs);
\item there is a set $Z \subseteq \{a, a{+}1, a{+}2, \ldots\}$ with $a \in Z$
and $\sum_{i \in Z} i! = 0$ in $\Zp$ (the sum converges since $\vpp(i!) \to \infty$);
\item there is an infinite chain $a = a_1 < a_2 < \cdots$ (the sorted elements
of $Z$) with $m_k := \vpp\bigl(\sum_{i \le k} a_i!\bigr) \to \infty$
--- the third sub-question of Erd\H{o}s \#404, answered \textbf{YES} for this $(a,p)$.
\end{enumerate}
\end{theorem}
\begin{proof}
(1)$\Rightarrow$(2): consider the tree whose nodes at depth $k$ are the
level-$k$ witnesses ($B \subseteq [a{+}1, i_{\max}(k)]$ with
$\sum_{j\in B}P_j \equiv -1 \pmod{p^k}$), with an edge from a depth-$(k{+}1)$
node $B_{k+1}$ to the depth-$k$ node $B_k = B_{k+1} \cap [a{+}1, i_{\max}(k)]$
(a level-$k$ witness, since terms with $\vpp(P_i) \ge k$ vanish mod $p^k$).
Every level is nonempty by (1) and each level is finite, so by K\"onig's
lemma there is an infinite branch $B_1, B_2, \ldots$ with
$B_k \subseteq B_{k+1}$; let $Z' := \bigcup_k B_k$ and $Z := Z' \cup \{a\}$.
For every $k$: $\sum_{i \in Z'}P_i \equiv \sum_{i\in B_k}P_i \equiv -1
\pmod{p^k}$ (terms beyond $i_{\max}(k)$ are $\equiv 0$), hence
$\sum_{i\in Z'}P_i = -1$ in $\Zp$ and $\sum_{i\in Z}i! =
a!\,(1 + \sum_{i \in Z'}P_i) = 0$ in $\Zp$.
(2)$\Rightarrow$(3): the partial sums satisfy
$\sum_{i \in Z, i \le N} i! = -\sum_{i \in Z, i > N} i!$, and the right side
is divisible by $\min_{i > N} \vpp(i!) \to \infty$.
(3)$\Rightarrow$(1): each $m_k$ certifies $-1 \in R_{m_k}$.
\end{proof}

\section{The $p \mid a{+}1$ class (Theorem B, kernel-green)}

\begin{theorem}\label{thm:B}
If $p \mid (a+1)$ then $P_i \equiv 0 \pmod p$ for every $i \ge a{+}1$, hence
$1 + \sum_{i \in B} P_i \equiv 1 \pmod p$ and $\vpp(1+\sum P_i) = 0$ for
\emph{every} $B$.  Therefore
\[ \boxed{f(a,p) = \vpp(a!) \quad \text{for all } p \mid a{+}1.} \]
\end{theorem}
\begin{proof}
$P_{a+1} = a{+}1 \equiv 0$ and $P_i = P_{a+1} \cdot (a{+}2)\cdots i$.
(Formalized for $p = 2$, all odd $a$: \texttt{erdos404\_odd\_a}.)
\end{proof}

\section{$f(4,2) = 6$ exactly (Theorem C, kernel-green)}

\begin{theorem}\label{thm:C}
$f(4,2) = 6$: the witness $4! + 5! + 7! = 5184 = 2^6 \cdot 81$, and every
sum $4! + \sum_B i!$ has $v_2 \le 6$.
\end{theorem}
\begin{proof}
Here $P_5 = 5, P_6 = 30, P_7 = 210$, and $P_i \equiv 0 \pmod{16}$ for all
$i \ge 8$ (the staircase acquires the factor $16 = 2^4$ at $i = 8$:
$P_8 = 1680 = 16\cdot 105$, and later terms only accumulate factors).
Modulo 16 the eight sub-sums $1 + \sum_{i \in S} P_i$ over
$S \subseteq \{5,6,7\}$ take values
$\{1, 6, 15, 3\} \cup \{4, 8, 1, 6\}$-class $\subseteq \{1,3,4,6,8,15\}$:
none $\equiv 0$.  Hence $\vpp(1+\sum_B P) \le 3$ for every $B$, and
$v_2 \le v_2(4!) + 3 = 6$.  The witness attains 6.
(Formalized: \texttt{erdos404\_f42}.)
\end{proof}

\section{The block-identity structure theory}

Consecutive factorial blocks factor exactly (kernel-green, first two):
\[
  [u, u{+}1] := u! + (u{+}1)! = u!\,(u+2),
\]
\[
  [u, u{+}2] = u!\,\bigl(1 + (u{+}1) + (u{+}1)(u{+}2)\bigr) = u!\,(u+2)^2,
\]
\[
  [u, u{+}3] = u!\,(u+2)\bigl(u^2 + 5u + 5\bigr),
\]
the first by $u! + (u{+}1)! = u!(1 + u + 1)$; the second because
$1 + (u{+}1) + (u{+}1)(u{+}2) = (u{+}2)^2$; the third because
$1 + (u{+}1) + (u{+}1)(u{+}2) + (u{+}1)(u{+}2)(u{+}3) = (u+2)(u^2+5u+5)$.
The \textbf{pair valuation law}
\[ v_2[u, u{+}1] = v_2(u!) + v_2(u+2) \]
is precisely the Kill-\#4 kernel law $v_2(\sum_{a\ge u}a!) = v_2(u!) + v_2(u+2)$
(the tail's valuation is realized by its first two terms).
These identities structure the gadget supply of Lemma~\ref{lem:lift}:
singles have valuations $v_2(i!)$ (the set $V = \{i - s_2(i)\}$ with its
digit-sum gaps), pairs fill $v_2(u!) + v_2(u+2)$, square-triples fill
$v_2(u!) + 2\,v_2(u+2)$, and the quadratic quads fill
$v_2(u!) + v_2(u+2) + v_2(u^2+5u+5)$.

\section{The computational classification (exact engines)}

\texttt{classifier.py} computes $g(a,p)$ exactly whenever the first
obstruction appears within the bitset budget; results
($a \le 40$, primes $\le 61$; $0^*$ = the Theorem-B class, proven):
\begin{itemize}
\item \textbf{p=2:} bounded: $g(4,2)=3,\ g(8,2)=6,\ g(10,2)=6,\ g(12,2)=3,\
g(20,2)=3,\ g(22,2)=11,\ g(26,2)=6,\ g(28,2)=3,\ g(32,2)=6,\ g(36,2)=3,\
g(40,2)=7$ (and $g=0$ for all odd $a$, Theorem B).
Unbounded candidates (no obstruction to depth $2^{28}$):
$a \in \{2, 6, 14, 16, 18, 24, 30, 34, 38, \ldots\}$.
\item \textbf{p=3:} bounded: $g(3,3)=2,\ g(6,3)=1,\ g(15,3)=1,\ g(21,3)=3,\
g(22,3)=3,\ g(24,3)=1,\ g(30,3)=2,\ g(33,3)=1,\ g(39,3)=7$; candidates at
depth $3^{13}$: most $a \equiv 1 \pmod 3$.
\item \textbf{p=5,7:} bounded: $g(2,5)=g(7,5)=g(12,5)=g(17,5)=g(22,5)=g(27,5)=g(32,5)=0$
(the window below 5 is too thin: $\{3!, 4!\} = \{6, 24\} \equiv \{1, 4\} \bmod 5$
has subset sums $\{0,1,4,0\}$, missing $-1$); $g(2,7)=g(7,7)=g(12,7)=g(13,7)=0$; etc.
\end{itemize}

\section{Conjecture D and the (2,2) chain}

\begin{conjecture}\label{conj:D}
$g(2,2) = \infty$ (and $g(6,2) = \infty$).  Equivalently: the set
$Z \subseteq \{2, 3, \ldots\}$ with $\sum_{i \in Z} i! = 0$ in $\mathbb{Z}_2$
exists, and the infinite chain of Erd\H{o}s's third sub-question starts
$2, 3, 5, 6, 7, 8, 9, 11, \ldots$
\end{conjecture}

\textbf{Evidence.} (i) No obstruction to depth $2^{28}$ (28 consecutive
levels, complete DP).  (ii) The explicit nested (K\"onig) chain, verified by
exact big-integer arithmetic (each line: $v_2(2! + \sum i!)$):
\[
\{3\}{:}3,\quad \{3,5\}{:}7,\quad \{3,5,8\}{:}9,\quad
\{3,5,6,7,8,9\}{:}9,\quad \{3,5,6,7,11,13,14\}{:}13,
\]
\[
\{3,5,6,7,11,12,15,16,19\}{:}18,\quad
\{3,5,6,7,11,12,15,17,19,21\}{:}23.
\]
(iii) The structure theory: the level-$k$ lift needs a fresh exact-valuation-$k$
gadget; the supply census (\texttt{zoo.py}, \texttt{gadget\_plan.json}) shows
singles/pairs/square-triples cover all $k$ except the digit-sum gaps
$k \in \{5, 6, 9, 12, 13, 14, \ldots\}$, and the DP proves those levels are
nonetheless reachable (by non-block subsets, e.g.\ $\{10,11,13\}$ has exact
valuation 9).  The proof frontier is a closed-form gadget family covering
every gap $k$ --- the remaining obstacle to a complete unconditional proof.

\section{What remains open toward the full \#404 classification}

\begin{enumerate}
\item The general classification: characterize $\{(a,p) : g(a,p) = \infty\}$.
The data suggests a dichotomy with genuinely unbounded cases (Conjecture D)
and a rich bounded set with exact small $g$-values.
\item The behavior of $f(a,p)$ on the bounded set (exact $g$-certificates
exist for every entry in the table; a closed form is not yet isolated).
\item The gap-valuation gadget families (Lemma~\ref{lem:lift} supply).
\end{enumerate}

\section{Prior art}
The site lists the problem as open (last edit 2025-12-01), tags number
theory / factorials, no proof expositions.  Kill \#4 of this campaign
(\texttt{problem404-padic-kernel/KernelLaw.lean}) is the full-tail case
$p = 2$; the present package generalizes to arbitrary index sets and primes.

\end{document}
